\documentclass[aspectratio=169,11pt]{beamer}
\usepackage{../phanes-beamer}
\renewcommand{\ModuleID}{03}
\title{Agentic software development}
\subtitle{Planner, researcher, implementer, verifier, tests, documentation, and the user}
\author{PHANES Initiative}
\institute{The University of Texas at Austin}
\date{September 12, 2026}
\begin{document}
\begin{frame}[plain]
\titlepage
\end{frame}
\begin{frame}[t,fragile]{One loop, several clearly bounded roles}
\fontsize{11.5}{14}\selectfont
\begin{itemize}\setlength{\itemsep}{2mm}
\item Planner: convert the user's objective into a small sequence with acceptance conditions and unresolved decisions.
\item Researcher: inspect the repository and primary documentation; return facts with paths, versions and sources.
\item Implementer: make the agreed code change. Verifier: challenge the diff and the evidence independently.
\item Test writer: construct requirement-driven checks. Documentation writer: describe the behavior that actually passed.
\item The user owns scope, consequential assumptions, resource limits, and final acceptance.
\end{itemize}
\end{frame}
\begin{frame}[t,fragile]{The development loop includes feedback}
\begin{center}
\begin{tikzpicture}[node distance=5mm and 9mm, every node/.style={draw=UTorange,rounded corners,align=center,font=\small,minimum height=7mm}, >={Stealth}]
\node (user) {User: scope and acceptance};
\node[below left=of user] (plan) {Planner};
\node[below right=of user] (research) {Researcher};
\node[below=15mm of user] (impl) {Implementer};
\node[below left=of impl] (verify) {Verifier};
\node[below right=of impl] (test) {Test writer};
\node[below=16mm of impl] (docs) {Documentation + user review};
\draw[->] (user)--(plan); \draw[->] (user)--(research);
\draw[->] (plan)--(impl); \draw[->] (research)--(impl);
\draw[->] (impl)--(verify); \draw[->] (impl)--(test);
\draw[->] (verify)--(docs); \draw[->] (test)--(docs);
\draw[->] (verify.west) to[bend left=45] (impl.west);
\end{tikzpicture}\end{center}
\end{frame}
\begin{frame}[t,fragile]{Define a handoff artifact for every role}
\fontsize{10.5}{12.5}\selectfont\setlength{\tabcolsep}{4pt}\renewcommand{\arraystretch}{1.15}
\begin{tabular}{>{\raggedright\arraybackslash}p{\dimexpr 0.2\linewidth-3.2pt\relax}>{\raggedright\arraybackslash}p{\dimexpr 0.8\linewidth-12.8pt\relax}}
\toprule
\textbf{Role} & \textbf{Handoff that another role can inspect} \\ \midrule
Planner & Task decomposition, files in scope, assumptions, risks, acceptance tests and stop conditions. \\[2pt]
Researcher & Repository facts and API references; source/version, relevant excerpt, uncertainty and proposed interface. \\[2pt]
Implementer & Minimal patch, explanation of changed behavior, commands run and unresolved failures. \\[2pt]
Verifier & Findings tied to requirements, reproducible failures, evidence gaps and a disposition. \\[2pt]
Test / docs writers & Independent tests and expected outcomes; then runnable instructions, limitations and examples. \\[2pt]
\bottomrule\end{tabular}\par\vspace{2mm}
\end{frame}
\begin{frame}[t,fragile]{Start with a task contract}
\fontsize{11.5}{14}\selectfont
\begin{itemize}\setlength{\itemsep}{2mm}
\item Example feature: export validated calculation results to CSV while preserving case identifiers and SI units.
\item Scope: edit the exporter and tests; do not change the physics calculator or accepted reference values.
\item Definition of done: correct columns, deterministic ordering, invalid-record rejection, and a documented CLI example.
\item Budget: two implementation/verification cycles before returning a concrete blocker; no network installs or publishing without a new decision.
\end{itemize}
\end{frame}
\begin{frame}[t,fragile]{What the user should ask the planner}
\begin{lstlisting}
Plan a CSV exporter for the existing result schema.
Inspect the repository first. Identify the files to change,
schema assumptions, invalid-input behavior, and acceptance tests.
Do not modify the calculator or relax existing tests.
Ask only about decisions that change the contract.
\end{lstlisting}
\fontsize{10.5}{12.5}\selectfont
\begin{itemize}\setlength{\itemsep}{2mm}
\item A useful response names concrete files and unresolved choices, not a generic software-development checklist.
\item Approve the plan's scope and test criteria; correct any mistaken physical or project assumption now.
\end{itemize}
\end{frame}
\begin{frame}[t,fragile]{Research should resolve implementation uncertainty}
\fontsize{11.5}{14}\selectfont
\begin{itemize}\setlength{\itemsep}{2mm}
\item Inspect existing schemas, tests, dependencies and naming conventions before searching for a new library.
\item For an unfamiliar API, consult the exact version's primary documentation and identify the required call signature.
\item Return the chosen approach with supporting paths or URLs and a short account of rejected alternatives.
\item Mark unknowns explicitly. Do not turn a model's recollection into a fact or introduce dependencies without a reason.
\end{itemize}
\end{frame}
\begin{frame}[t,fragile]{The implementer works inside the accepted scope}
\fontsize{11.5}{14}\selectfont
\begin{itemize}\setlength{\itemsep}{2mm}
\item Change the smallest coherent set of files that satisfies the contract.
\item Preserve unrelated user edits and existing public interfaces unless the agreed change requires otherwise.
\item Run the relevant checks and retain exact commands, exit codes and output artifacts.
\item When an assumption fails, stop the dependent work and return the specific fact that needs resolution.
\end{itemize}
\end{frame}
\begin{frame}[t,fragile]{The verifier must be able to disagree}
\fontsize{11.5}{14}\selectfont
\begin{itemize}\setlength{\itemsep}{2mm}
\item Start from the requirement and the actual diff, not just the implementer's summary.
\item Inspect edge cases: missing fields, invalid units, stale outputs, partial writes and changed assumptions.
\item Reproduce a representative success and at least one failure path where execution access is authorized.
\item Return a specific finding or a scoped acceptance statement, with the evidence and remaining limits.
\end{itemize}
\end{frame}
\begin{frame}[t,fragile]{Test design comes from requirements}
\fontsize{10.5}{12.5}\selectfont\setlength{\tabcolsep}{4pt}\renewcommand{\arraystretch}{1.15}
\begin{tabular}{>{\raggedright\arraybackslash}p{\dimexpr 0.28\linewidth-4.48pt\relax}>{\raggedright\arraybackslash}p{\dimexpr 0.72\linewidth-11.52pt\relax}}
\toprule
\textbf{Requirement} & \textbf{Test that can distinguish correct from incorrect} \\ \midrule
Preserve identity & Two equal values with different case IDs must remain distinct rows. \\[2pt]
Reject invalid units & A record labeled degrees C where K is required must fail explicitly. \\[2pt]
Avoid partial success & One malformed record must not leave an apparently complete output file. \\[2pt]
Stable output & Different input ordering produces the declared canonical row order. \\[2pt]
No stale result & Changed input cannot reuse an output bound to an earlier input hash. \\[2pt]
\bottomrule\end{tabular}\par\vspace{2mm}
\end{frame}
\begin{frame}[t,fragile]{Documentation follows observed behavior}
\fontsize{11.5}{14}\selectfont
\begin{itemize}\setlength{\itemsep}{2mm}
\item Describe prerequisites, configuration variables, command syntax, input/output schemas and one checked example.
\item State known limitations and failure behavior, including what the command will not overwrite.
\item Link acceptance evidence and identify any step that was not executed in the current environment.
\item When implementation changes the interface, update both the example and its associated test.
\end{itemize}
\end{frame}
\begin{frame}[t,fragile]{The user has four main interaction points}
\fontsize{10.5}{12.5}\selectfont\setlength{\tabcolsep}{4pt}\renewcommand{\arraystretch}{1.15}
\begin{tabular}{>{\raggedright\arraybackslash}p{\dimexpr 0.24\linewidth-3.84pt\relax}>{\raggedright\arraybackslash}p{\dimexpr 0.76\linewidth-12.16pt\relax}}
\toprule
\textbf{When} & \textbf{Useful user action} \\ \midrule
Before planning & State objective, context, accepted inputs, constraints and required output. \\[2pt]
At plan review & Resolve material ambiguity; accept files in scope, test criteria and resource budget. \\[2pt]
At exceptions & Decide changed physics, new dependency, larger compute job or external disclosure. \\[2pt]
At handoff & Review diff, test evidence, unresolved findings, runnable docs and final result. \\[2pt]
\bottomrule\end{tabular}\par\vspace{2mm}
\end{frame}
\begin{frame}[t,fragile]{A verifier reports a real defect}
\fontsize{11.5}{14}\selectfont
\begin{itemize}\setlength{\itemsep}{2mm}
\item The exporter labels every row "K", but one input actually stores degrees C.
\item Verifier handoff: exact input file, observed row, violated unit contract and a reproducing command.
\item Implementer response: reject the malformed record or implement the explicitly agreed conversion.
\item Test writer adds the distinguishing case; docs describe the accepted unit contract. The user decides whether conversion changes scope.
\end{itemize}
\end{frame}
\begin{frame}[t,fragile]{Permissions should follow role responsibilities}
\fontsize{10.5}{12.5}\selectfont\setlength{\tabcolsep}{4pt}\renewcommand{\arraystretch}{1.15}
\begin{tabular}{>{\raggedright\arraybackslash}p{\dimexpr 0.23\linewidth-3.68pt\relax}>{\raggedright\arraybackslash}p{\dimexpr 0.77\linewidth-12.32pt\relax}}
\toprule
\textbf{Role} & \textbf{Default classroom authority} \\ \midrule
Planner / researcher & Read approved project files; no edits or shell execution by default. Research sources may be preloaded. \\[2pt]
Implementer & Read project files; request scoped edits and the relevant local test commands. \\[2pt]
Verifier & Read diff and evidence; request independent tests; no silent code repair. \\[2pt]
Test writer & Request edits under tests/ and approved test runs; no changes to acceptance requirements. \\[2pt]
Docs writer & Request edits under docs/ and checked examples; no publishing authority. \\[2pt]
\bottomrule\end{tabular}\par\vspace{2mm}
\end{frame}
\begin{frame}[t,fragile]{OpenCode permission syntax is explicit}
\begin{lstlisting}
"permission": {
  "*": "ask",
  "read": {"*": "allow", ".env": "deny", "*.env": "deny"},
  "edit": "ask",
  "bash": "ask",
  "webfetch": "deny"
}
\end{lstlisting}
\fontsize{10.5}{12.5}\selectfont
\begin{itemize}\setlength{\itemsep}{2mm}
\item allow runs automatically; ask requires an approval; deny blocks the matching action.
\item Specific rule order matters. Inspect merged project, global and agent settings.
\item Shell execution can read files or use the network; file-tool rules alone are not OS enforcement.
\end{itemize}
\vfill{\fontsize{6.5}{8}\selectfont\color{UTgray} Sources: ocpermissions, occonfig. See references and instructor notes.\par}
\end{frame}
\begin{frame}[t,fragile]{Define a read-only planning role}
\begin{lstlisting}
"agent": {"course-planner": {
  "mode": "primary",
  "description": "Plan bounded engineering software changes",
  "prompt": "Inspect the contract; propose files, tests and decisions.",
  "permission": {"edit": "deny", "bash": "deny",
                 "webfetch": "deny", "task": "deny"}
}}
\end{lstlisting}
\fontsize{10.5}{12.5}\selectfont
\begin{itemize}\setlength{\itemsep}{2mm}
\item This example deliberately starts without delegation authority.
\item Add role-specific delegation only when the invoked subagent's own permissions have been reviewed.
\end{itemize}
\vfill{\fontsize{6.5}{8}\selectfont\color{UTgray} Sources: ocagents. See references and instructor notes.\par}
\end{frame}
\begin{frame}[t,fragile]{Allow only the intended delegated roles}
\begin{lstlisting}
"permission": {"task": {
  "*": "deny",
  "course-researcher": "allow",
  "course-implementer": "ask",
  "course-verifier": "allow",
  "course-test-writer": "ask",
  "course-doc-writer": "ask"
}}
\end{lstlisting}
\fontsize{10.5}{12.5}\selectfont
\begin{itemize}\setlength{\itemsep}{2mm}
\item This task rule belongs to the proposed coordinator, not the read-only example on the prior slide.
\item Invocation permission does not define the invoked role's file and execution authority.
\item OpenCode also supports direct user @ invocation; role access must stand on its own.
\end{itemize}
\vfill{\fontsize{6.5}{8}\selectfont\color{UTgray} Sources: ocagents. See references and instructor notes.\par}
\end{frame}
\begin{frame}[t,fragile]{A permitted command can still have broad effects}
\fontsize{11.5}{14}\selectfont
\begin{itemize}\setlength{\itemsep}{2mm}
\item "Run pytest" executes repository code; test fixtures may write files, launch subprocesses or contact services.
\item "Run Python" is a general execution capability, not a calculator-only permission.
\item Use an unprivileged course account, reviewed dependencies, bounded working directories and appropriate network controls.
\item Keep credentials outside the agent-visible task files; review commands and environment exposure before granting execution.
\end{itemize}
\end{frame}
\begin{frame}[t,fragile]{Use Git boundaries for coherent changes}
\fontsize{11.5}{14}\selectfont
\begin{itemize}\setlength{\itemsep}{2mm}
\item Begin with the current status and distinguish existing user edits from the requested task.
\item Assign one owner to each edited file or use separate branches/worktrees when concurrent edits are useful.
\item Review a coherent diff after each implementation cycle; keep failed experiments out of accepted artifacts.
\item Merging, publishing, and changing shared infrastructure require the applicable user/project authority.
\end{itemize}
\end{frame}
\begin{frame}[t,fragile]{Agent, skill, and tool remain distinct}
\fontsize{10.5}{12.5}\selectfont\setlength{\tabcolsep}{4pt}\renewcommand{\arraystretch}{1.15}
\begin{tabular}{>{\raggedright\arraybackslash}p{\dimexpr 0.23\linewidth-3.68pt\relax}>{\raggedright\arraybackslash}p{\dimexpr 0.77\linewidth-12.32pt\relax}}
\toprule
\textbf{Object} & \textbf{Example in this development loop} \\ \midrule
Agent role & Verifier with a review prompt and restricted edit authority. \\[2pt]
Skill & Repeatable CSV-contract review procedure with required checks and stop conditions. \\[2pt]
Tool & Read a file, inspect a diff, execute a test or parse a result artifact. \\[2pt]
Evidence & The actual diff, process status, test report and saved output under review. \\[2pt]
\bottomrule\end{tabular}\par\vspace{2mm}
\end{frame}
\begin{frame}[t,fragile]{A skill expresses a repeatable review procedure}
\begin{lstlisting}
# .opencode/skills/result-review/SKILL.md
---
name: result-review
description: Check result identity, units and execution provenance.
---
Require the input manifest, tool status and output artifact.
Compare identifiers and units against the accepted contract.
Reject missing, stale or nonfinite results.
Report evidence paths and unresolved findings.
\end{lstlisting}
\vfill{\fontsize{6.5}{8}\selectfont\color{UTgray} Sources: ocskills. See references and instructor notes.\par}
\end{frame}
\begin{frame}[t,fragile]{Bound retries and return an actionable blocker}
\fontsize{11.5}{14}\selectfont
\begin{itemize}\setlength{\itemsep}{2mm}
\item Retry a transient process failure only within the declared attempt/time budget.
\item Do not retry an invalid physical assumption as if it were a temporary network problem.
\item After repeated verification failure, report the smallest reproducer, suspected cause, failed remedies and remaining decision.
\item Resume from that evidence after the user resolves the blocker; do not restart the project or conceal failed trials.
\end{itemize}
\end{frame}
\begin{frame}[t,fragile]{OAK maps the bounded loop to named agents}
\fontsize{10.5}{12.5}\selectfont\setlength{\tabcolsep}{4pt}\renewcommand{\arraystretch}{1.05}
\begin{tabular}{>{\raggedright\arraybackslash}p{\dimexpr 0.23\linewidth-3.68pt\relax}>{\raggedright\arraybackslash}p{\dimexpr 0.77\linewidth-12.32pt\relax}}
\toprule
\textbf{Course role} & \textbf{OAK agent} \\ \midrule
Coordinator / planner & \texttt{lead}: bounded router; simple low-risk work goes directly to \texttt{developer}. \\[2pt]
Researcher & \texttt{researcher}: code, documentation, API, alternative and risk investigation. \\[2pt]
Planner & \texttt{specifier}: tasks, acceptance criteria, scope boundaries and validation plan. \\[2pt]
Implementer & \texttt{developer}: implementation, direct-mode changes and focused validation. \\[2pt]
Verifier & \texttt{reviewer}: final read-only review against scope and the actual diff. \\[2pt]
\bottomrule\end{tabular}\par\vspace{1mm}
\fontsize{10.5}{12.5}\selectfont
\begin{itemize}\setlength{\itemsep}{2mm}
\item OAK (Apache-2.0) is an open-source OpenCode harness that keeps the loop local, inspectable, and reproducible.
\item Its shipped default permissions match the course defaults: reads allowed, edits and bash ask, external directories denied.
\end{itemize}
\vfill{\fontsize{6.5}{8}\selectfont\color{UTgray} Sources: oak. See references and instructor notes.\par}
\end{frame}
\begin{frame}[t,fragile]{OAK's bounded loops are a concrete safety contract}
\fontsize{11.5}{14}\selectfont
\begin{itemize}\setlength{\itemsep}{2mm}
\item \texttt{/loop}: writes require an approved task contract; at most three iterations per invocation within a 1--6-task budget; resumable from durable state.
\item \texttt{/autonomous}: local-checkout-only, deterministic validation each iteration, stops when the independent \texttt{reviewer} approves.
\item The \texttt{/autonomous} contract prohibits network access, commits, pushes, merges, deployments and publication.
\item State uses a schema-versioned JSON snapshot, append-only JSONL history, exclusive lock, contract hashing and idempotent action IDs.
\item The harness itself is mechanically validated: agent and command contract checks plus 1{,}000+ automated tests.
\end{itemize}
\fontsize{10.5}{12.5}\selectfont
\begin{itemize}\setlength{\itemsep}{1mm}
\item A second reference implementation of the same principles: OAK targets software engineering with any provider, while PHANES applies them to private-endpoint reactor calculations.
\end{itemize}
\vfill{\fontsize{6.5}{8}\selectfont\color{UTgray} Sources: oak. See references and instructor notes.\par}
\end{frame}
\begin{frame}[t,fragile]{A useful completion message closes the contract}
\fontsize{11.5}{14}\selectfont
\begin{itemize}\setlength{\itemsep}{2mm}
\item State the changed behavior and point to the actual code or artifact.
\item Summarize the relevant checks, including failed or unexecuted ones.
\item List unresolved limitations and decisions that affect use.
\item Give a runnable handoff command and the evidence needed for user acceptance.
\end{itemize}
\end{frame}
\begin{frame}[t,fragile]{Lab: demonstrate the complete development loop}
\fontsize{11.5}{14}\selectfont
\begin{itemize}\setlength{\itemsep}{2mm}
\item Implement the CSV feature using the proposed role sequence, with explicit task and handoff records.
\item Retain one verifier-found defect and show its test, repair and recheck.
\item Document role permissions and demonstrate that a disallowed edit is blocked in the reviewed setup.
\item Submit the final diff, acceptance evidence, usage documentation and the user decisions that changed the work.
\end{itemize}
\end{frame}
\begin{frame}[t,fragile]{Sample submission: the development-loop lab}
\fontsize{11.5}{14}\selectfont
\begin{itemize}\setlength{\itemsep}{2mm}
\item Task and handoff records: each role records its inputs, artifacts, and acceptance criteria; the planner issues the task spec, the implementer returns the diff, the verifier returns the signed check list.
\item Retained defect: the verifier's test found a column written without its unit header; the submission shows the failing test, the repair, the recheck passing, and the diff between the two runs.
\item Blocked edit: the permission setting denies writes outside the lab directory; the demonstration shows the denial message and confirms that no file changed.
\item Final package: the complete diff, passing test output, usage documentation, and a short list of user decisions (scope, naming, stopping criteria) that changed the work.
\end{itemize}
\end{frame}
\begin{frame}[t,fragile]{Exit ticket: bounded roles in practice}
\fontsize{11.5}{14}\selectfont
\begin{itemize}\setlength{\itemsep}{2mm}
\item Why must the verifier role be denied permission to edit code directly?
\item What mechanism prevents the agent loop from cycling indefinitely on a failing test?
\item Why does a tool allowlist in the harness fail to guarantee operating-system security?
\item When must an unexpected observation escalate to the user rather than undergoing autonomous retry?
\end{itemize}
\end{frame}
\begin{frame}[t,fragile]{A complete answer: bounded roles}
\fontsize{11.5}{14}\selectfont
\begin{itemize}\setlength{\itemsep}{2mm}
\item Independent challenge: allowing the verifier to repair code merges advocacy with verification, risking silent test relaxation or scope creep.
\item Bounded retry: the planner or harness enforces an explicit iteration budget (e.g. 2--3 attempts) before converting a persistent failure into an actionable blocker.
\item Shell execution escape: a permitted shell command (such as pytest or python) can execute arbitrary system calls and access the network regardless of file-tool rules.
\item Escalation threshold: escalate to the user whenever an assumption alters physical inputs, interfaces, compute budgets, or regulatory information boundaries.
\end{itemize}
\end{frame}
\begin{frame}[t]{References 1}
\fontsize{9}{12}\selectfont\begin{itemize}\setlength{\itemsep}{3mm}
\item \textbf{curriculum}: User-supplied PHANES curriculum: mod/01 through mod/08/README.md.
\item \textbf{colors}: \href{https://umac.utexas.edu/brand-center/colors/}{umac.utexas.edu/brand-center/colors}
\item \textbf{ocpermissions}: \href{https://opencode.ai/docs/permissions/}{opencode.ai/docs/permissions}
\item \textbf{occonfig}: \href{https://opencode.ai/docs/config/}{opencode.ai/docs/config}
\item \textbf{ocagents}: \href{https://opencode.ai/docs/agents/}{opencode.ai/docs/agents}
\item \textbf{ocskills}: \href{https://opencode.ai/docs/skills/}{opencode.ai/docs/skills}
\item \textbf{oak}: \href{https://github.com/jcarlosrodicio/opencode-agent-orchestration-kit}{github.com/jcarlosrodicio/opencode-agent-orchestration-kit}
\end{itemize}\end{frame}
\end{document}